% SPDX-License-Identifier: GPL-3.0-or-later
% Copyright (C) 2026 Alessandro Greco
\section*{Глава \ref{cap:07-canali}: каналы и протокол связи}

\begin{soluzione}{es:07-ruoli}
Клиент всегда сам инициирует соединения: с входным узлом (guard) и ретрансляторами, предоставляющими сведения системы каталогов. Никому не требуется подключаться к клиенту, а прослушивание порта лишь облегчило бы его идентификацию. Ретранслятор инициирует каналы главным образом тогда, когда получает запрос \texttt{EXTEND2} и должен связаться со следующим переходом цепи, либо когда, как любой узел, создаёт собственные цепи для загрузки документов каталогов или проверки своей доступности.
\end{soluzione}

\begin{soluzione}{es:07-scadenza}
$\texttt{0x000772F7} = 488183$ часа, что соответствует 9 сентября 2025 года, 23:00 UTC. 32-битный беззнаковый счётчик секунд переполнился бы в 2106 году. Подсчёт часов вместо секунд даёт запас в сотни тысяч лет, а точности до часа достаточно для сертификатов, действующих дни или недели. Одно примечание спецификации требует осторожности: \texttt{cert-spec} утверждает, что четырёх байт достаточно до 10136 года. Это ошибка, явно отмеченная в \proposta[Предложении]{338}; она возникла из-за интерпретации счётчика как измеряющего \emph{минуты}, поскольку $2^{32}$ минут — примерно 8\,166 лет.
\end{soluzione}

\begin{soluzione}{es:07-authchallenge}
$38 = 32 + 2 + 2 \times 2$ описывает проверочный вызов с двумя методами. Если N\_Methods равно 2, тело должно содержать $32 + 2 + 4 = 38$ байт. Тело ячейки длиной 36 байт объявляет больше методов, чем фактически содержит, и является некорректным, поэтому должно быть отклонено. Спецификация разрешает игнорировать \emph{лишние} байты, но не отсутствующие.
\end{soluzione}

\begin{soluzione}{es:07-impostore}
Сравнение ключа, сертифицированного сертификатом типа 5, с отпечатком SHA-256 TLS-сертификата атакующего завершится неудачей: он отличается от сертификата настоящего ретранслятора. Если атакующий предъявит тот же подлинный TLS-сертификат, он не сможет завершить рукопожатие TLS без соответствующего закрытого ключа.
\end{soluzione}

\begin{soluzione}{es:07-netinfo}
\texttt{695661c0} (1 января 2026 года, 12:00:00 UTC); \texttt{04 04 c6336407} (адрес клиента, 198.51.100.7); \texttt{01}; и \texttt{04 04 cb007105} (адрес ретранслятора, 203.0.113.5), затем заполнение нулевыми байтами до общей длины 509 байт.
\end{soluzione}

\begin{soluzione}{es:07-client-auth}
Во-первых, входной узел видел бы один и тот же идентификатор, когда пользователь переключается между домашней, рабочей и мобильной сетями, связывая сеансы и IP-адреса во всех этих средах. Во-вторых, идентификатор также передавался бы ретрансляторам, с которых клиент загружает документы каталогов. Сотрудничающие ретрансляторы либо атакующий, наблюдающий их, могли бы связывать активность пользователя между разными ретрансляторами. В обоих случаях постоянный идентификатор разрушал бы предполагаемую Tor несвязываемость.
\end{soluzione}

\begin{soluzione}{es:07-rsa}
Общий план: разобрать сертификат типа 2 библиотекой X.509, например \texttt{cryptography.x509}. Извлечь его открытый RSA-ключ, проверить самоподпись и убедиться, что показатель степени равен 65537. Затем разобрать перекрёстный сертификат типа 7 согласно формату. Убедиться, что сертифицированный Ed25519-ключ совпадает с уже проверенной идентичностью, и проверить RSA-подпись с PKCS\#1 v1.5 без OID алгоритма хэширования над SHA-256-хэшем предписанного префикса и полей. Отпечаток ретранслятора — дайджест SHA-1 DER-кодированного представления открытого ключа RSA в формате PKCS\#1.
\end{soluzione}

\begin{soluzione}{es:07-padding}
Нет. Трафик клиента \texttt{test008c} проходит через \emph{собственный} канал к его входному узлу, а не через наш канал. Наш канал никогда не использовался для цепи, поэтому \file{channelpadding\_decide\_to\_pad\_channel} решает не добавлять заполнение. Чтобы наблюдать заполнение канала, нужно передавать многоузловую цепь через собственный канал, как в \cref{cap:10-cifratura-relay}.
\end{soluzione}

\begin{soluzione}{es:07-mitm}
TLSSECRETS выводится из секретов TLS-сессии между $A$ и посредником, и подпись $A$ покрывает это значение. $B$ повторно вычисляет TLSSECRETS из секретов \emph{своей} сессии с посредником. Эти секреты различаются, поэтому проверка не проходит даже тогда, когда посредник честно пересылает каждую ячейку и SLOG с CLOG совпадают. Без TLSSECRETS и SCERT посредник мог бы представиться $B$ как $A$ и переслать подпись $A$, заставив $B$ ошибочно считать, что у него аутентифицированный канал с $A$.
\end{soluzione}

\section*{Глава \ref{cap:08-celle}: ячейки}

\begin{soluzione}{es:08-lunghezze}
\cella{CREATE2} — ячейка фиксированной длины: 512 байт при протоколе связи версии 3 с двухбайтовым CircID или 514 байт при версии 5 с четырёхбайтовым CircID. Ячейка \cella{VERSIONS} всегда использует двухбайтовый CircID, поэтому её длина равна $2 + 1 + 2 + 4 = 9$ байт.
\end{soluzione}

\begin{soluzione}{es:08-formato}
0 (\cella{PADDING}) — фиксированной длины; 7 (\cella{VERSIONS}) — переменной; 9 (\cella{RELAY\_EARLY}) — фиксированной; 12 (\cella{PADDING\_NEGOTIATE}) — фиксированной; 129 (\cella{CERTS}) — переменной. Команда 200 по правилу $\ge 128$ имела бы переменную длину, но она не определена и поэтому недопустима. Получатель может отбросить её и закрыть канал.
\end{soluzione}

\begin{soluzione}{es:08-circid}
Только \texttt{0x92345678}. CircID, равный нулю, зарезервирован для ячеек вне цепи. При протоколе связи версии $\ge 4$ сторона, инициирующая канал, здесь клиент, должна выбирать CircID со старшим битом, равным 1; у \texttt{0x12345678} старший бит (MSB) равен 0.
\end{soluzione}

\begin{soluzione}{es:08-hex}
Первая ячейка: CircID \texttt{00000000}, команда \texttt{0x81} = 129 (\cella{CERTS}), следовательно, длина переменная; длина \texttt{0003}, тело \texttt{aa bb cc}. Вторая ячейка: CircID \texttt{0x80000007}, команда \texttt{0x04} (\cella{DESTROY}), фиксированной длины, поэтому имеет тело 509 байт. Первый байт тела, \texttt{00}, означает причину NONE. Тело \texttt{aabbcc} само по себе не является корректным объектом CERTS: упражнение посвящено только формату обрамления ячеек.
\end{soluzione}

\begin{soluzione}{es:08-relay}
\texttt{02 0000 0005 00000000 012c}: команда DATA (2), Recognized 0, StreamID 5, Digest временно равен нулю, длина $300 = \texttt{0x012c}$. Оставшиеся $509 - 11 - 300 = 198$ байт — заполнение: первые четыре байта нулевые, остальные 194 случайные.
\end{soluzione}

\begin{soluzione}{es:08-recognized}
Для ячейки, всё ещё зашифрованной ещё одним слоем, Recognized равен нулю с вероятностью $2^{-16}$: в среднем примерно $10^6 / 65536 \approx 15$ таких ячеек на миллион. Одного Recognized недостаточно, поскольку возможны ложные совпадения. Только сравнение с накопительным дайджестом подтверждает, что ячейка адресована ретранслятору: для узла, не являющегося адресатом, случайное совпадение дайджеста потребовало бы дополнительной вероятности $2^{-32}$.
\end{soluzione}

\begin{soluzione}{es:08-destroy}
Предположим, атакующий контролирует входной и выходной ретрансляторы. Когда выходной ретранслятор хочет пометить цепь, он закрывает её редкой причиной, например \texttt{NOSUCHSERVICE}. Если промежуточный и входной ретрансляторы пересылали бы эту причину без изменений к клиенту, входной ретранслятор мог бы одновременно увидеть необычное значение в одной из своих цепей и связать известного ему клиента с назначением, известным выходному ретранслятору. Если все ретрансляторы передают и распространяют только NONE, байт причины не может служить таким сигналом. Корреляция по времени закрытия цепи всё ещё возможна, но этот канал намного более шумный.
\end{soluzione}

\begin{soluzione}{es:08-efficienza}
Традиционный формат: $\lceil 1048576/498 \rceil = 2106$ ячеек, занимающих $2106 \times 514 = \num{1082484}$ байт в канале, то есть избыточные данные составляют 3,2\%. CGO: $\lceil 1048576/488 \rceil = 2149$ ячеек, занимающих \num{1104586} байт, то есть избыточные данные составляют 5,3\%. Сообщения SENDME и открытия потока не учитываются.
\end{soluzione}

\begin{soluzione}{es:08-lab-fuzz}
Общий план: использовать Hypothesis для генерации списков кортежей, содержащих CircID, команду и тело ячейки. Команды выбирать из \texttt{CellCommand}, а тела ячеек фиксированной длины ограничивать 509 байтами. Объединить результаты \texttt{encode}, разбить полученную последовательность на фрагменты случайных длин и постепенно добавлять их в буфер. Многократно вызывать \texttt{decode}, извлекая ячейки до тех пор, пока не останется полной ячейки. Для фиксированных ячеек сравнивать тела с исходными, дополненными нулями до 509 байт.
\end{soluzione}

\begin{soluzione}{es:08-arti}
В Arti 2.6.0 кодек канала — \texttt{ChannelCodec} в \file{crates/tor-cell/src/chancell/codec.rs}, инициализируемый версией протокола связи через \texttt{ChannelCodec::new}. Его методы \file{write\_cell} и \file{decode\_cell} сериализуют и извлекают ячейки. Подобно нашему \texttt{decode}, \file{decode\_cell} обрабатывает неполные ячейки без ошибки, возвращая \texttt{Ok(None)}, пока не поступит достаточно байт, и удаляя данные из буфера только для полной ячейки. Arti, однако, работает с изменяемым буфером \texttt{BytesMut} и использует обобщённый тип сообщения канала (\texttt{M: ChanMsg}), позволяя сузить набор принимаемых сообщений. Ещё одно заметное отличие: \file{decode\_cell} читает команду по смещению 4 и предполагает фиксированные ячейки длиной 514 байт, то есть четырёхбайтовые CircID. Arti поддерживает только протоколы связи версий 4 и 5, как задаёт \file{LINK\_PROTOCOLS} в \file{crates/tor-proto/src/channel/handshake.rs}; начальная \cella{VERSIONS} с двухбайтовым CircID обрабатывается отдельно во время рукопожатия.
\end{soluzione}

\section*{Глава \ref{cap:09-circuiti}: построение цепи}

\begin{soluzione}{es:09-dimensioni}
ntor требует $20 + 32 + 32 = 84$ байта. ntor-v3 с трёхбайтовым сообщением требует $32 + 32 + 32 + 3 + 32 = 131$ байт. HDATA может содержать максимум $509 - 4 = 505$ байт, поэтому $CM$ может занимать не более $505 - 128 = 377$ байт.
\end{soluzione}

\begin{soluzione}{es:09-fast}
Секрет для \cella{CREATE\_FAST} равен $X \| Y$ и передаётся открытым текстом внутри канала. Для второго перехода значения $X$ и $Y$ прошли бы через цепь через первый ретранслятор. Этот ретранслятор знает ключи собственного слоя и смог бы прочитать значения, вывести ключи второго перехода и расшифровать всё, чем клиент обменивался с ним, уничтожая разделение между переходами.
\end{soluzione}

\begin{soluzione}{es:09-direzioni}
AES-CTR шифрует, выполняя XOR открытого текста с потоком ключа, зависящим только от ключа и счётчика. При одном и том же ключе в обоих направлениях ячейка клиента номер $n$ и ячейка ретранслятора номер $n$ использовали бы один и тот же поток ключа. Промежуточный ретранслятор, наблюдающий оба шифртекста, мог бы выполнить XOR и получить XOR двух открытых текстов, что нередко позволило бы восстановить оба сообщения.
\end{soluzione}

\begin{soluzione}{es:09-extend2}
Ретранслятор \texttt{test004r} использует ORPort 5104 = \texttt{0x13f0}: \texttt{03} означает три спецификатора соединения; спецификатор IPv4-адреса и порта равен \texttt{00 06 7f000001 13f0}. Полное сообщение содержит $1 + (2+6) + (2+20) + (2+32) = 65$ байт спецификаторов, плюс $2 + 2 = 4$ байта для HTYPE и HLEN и $32 \cdot 3 + 1 + 32 = 129$ байт HDATA. Итого 198 байт, что свободно помещается в 498-байтовую полезную ёмкость ретрансляционного сообщения.
\end{soluzione}

\begin{soluzione}{es:09-ds}
Без $\mathrm{EXP}(B, x)$ ретранслятор не аутентифицирован. Промежуточный ретранслятор, пересылающий рукопожатие, мог бы сгенерировать собственные $Y$ и $\mathit{AUTH}$ со своим $y$, выдать себя за следующий переход и читать последующий трафик, предназначенный ему: это \emph{man-in-the-middle}-атака. Без $\mathrm{EXP}(Y, x)$ теряется прямая секретность. Любой, кто позднее получит долгосрочный ключ onion $b$, сможет вычислить $\mathrm{EXP}(X, b)$ из записанных сообщений рукопожатия и вывести все сеансовые ключи.
\end{soluzione}

\begin{soluzione}{es:09-ver}
\texttt{VER} привязывает рукопожатие к его контексту, для Tor это \texttt{"circuit extend"}. Поскольку значение участвует во всех вычислениях ключевого материала, рукопожатие, созданное в одном контексте, не даёт ни ключей, ни значений $\mathit{AUTH}$, допустимых в другом. Без \texttt{VER} атакующий мог бы повторно использовать сообщение рукопожатия из другого протокола против ретранслятора Tor или наоборот, если обе стороны принимали одинаковые вычисления. Это \emph{межпротокольные} атаки.
\end{soluzione}

\begin{soluzione}{es:09-cc}
В лабораторной сети \texttt{test003r}, поддерживающий \texttt{Relay=2-6}, отклоняет запрос, состоящий только из \texttt{[02 06]}, ячейкой \cella{DESTROY} с причиной \texttt{PROTOCOL}. Комбинированный запрос с \file{CC\_FIELD\_REQUEST} он принимает и отвечает \file{CC\_FIELD\_RESPONSE}. Код tor включает CGO на переходе только вместе с управлением перегрузкой: \file{enable\_cgo = supports\_cgo \&\& use\_congestion\_control} в \file{src/core/or/extendinfo.c}. \proposta{359} объясняет, почему новый формат ретрансляционного сообщения, где StreamID имеется лишь у некоторых команд, несовместим с прежним механизмом управления потоком через SENDME на уровне потоков. Он требует управления потоком из \proposta[Предложения]{324}, то есть \texttt{FlowCtrl=2}.
\end{soluzione}

\begin{soluzione}{es:09-ntor-ref}
Это практическое упражнение. \file{ntor\_ref.py} содержит клиентскую и серверную функции. Использование фиксированных ключей вместо случайной генерации новых даёт воспроизводимые значения для сравнения с \file{ntor.py}.
\end{soluzione}

\begin{soluzione}{es:09-relay}
Общий план: \file{server\_part1} получает HDATA, проверяет \texttt{ID} и $B$, вычисляет $X_b = \mathrm{EXP}(X, b)$, выводит ключи первой фазы, проверяет MAC и расшифровывает $CM$. Затем \file{server\_part2} генерирует $(y, Y)$, вычисляет $\mathit{key\_seed}$ и $\mathit{verify}$, выводит ключ для $SM$ и вычисляет $\mathit{AUTH}$. Если $y$ зафиксировать на значении тестового вектора из предложения, ответ должен совпасть с \file{server\_handshake}.
\end{soluzione}

\section*{Глава \ref{cap:10-cifratura-relay}: луковица}

\begin{soluzione}{es:10-ordine}
Клиент обновляет накопительный дайджест прямого направления второго перехода, шифрует с Kf$_2$, а затем с Kf$_1$. Третий переход ничего не видит: второй переход распознаёт ячейку и не пересылает её дальше.
\end{soluzione}

\begin{soluzione}{es:10-contatore}
Каждая ячейка потребляет 509 байт потока ключа. После 1\,000 ячеек позиция в потоке ключа равна $509 \times 1000 = \num{509000}$ байт, то есть после AES-блока \num{31812}, поскольку блоки имеют длину 16 байт. Повторное использование ключа для новой цепи со счётчиком, начинающим снова с нуля, привело бы к повторному использованию потока ключа; XOR двух шифртекстов, зашифрованных им, раскрыл бы XOR исходных открытых текстов. Поэтому каждая цепь выводит свежие ключи через рукопожатие с эфемерными значениями.
\end{soluzione}

\begin{soluzione}{es:10-falso}
Накопительный дайджест ретранслятора включил бы ячейку, которую клиент не учитывал для этого перехода. После этого состояния дайджеста разошлись бы, и ни одна следующая ячейка не распознавалась бы, уничтожая цепь. Для каждой транзитной ячейки, не адресованной ретранслятору, проблема возникает с вероятностью $2^{-16}$: в среднем один раз на \num{65536} таких ячеек, что соответствует примерно 33\,MB трафика.
\end{soluzione}

\begin{soluzione}{es:10-tag}
Атакующему следует изменять байты, ожидаемый открытый текст которых известен или предсказуем, чтобы сообщник на выходном ретрансляторе мог обнаружить конкретное изменение: например команду ретрансляционного сообщения, \texttt{02} для \texttt{DATA}, StreamID либо четыре нулевых байта в начале заполнения. До проверки дайджеста сотрудничающий выходной ретранслятор смотрит, соответствует ли значение ожидаемой модификации. Честный выходной ретранслятор видит лишь ячейку с неверным дайджестом и уничтожает цепь.
\end{soluzione}

\begin{soluzione}{es:10-cgo-tag}
При CGO после снятия слоя шифрования выходного ретранслятора изменённая ячейка даёт псевдослучайные тег и содержимое. Узнаваемой метки нет: изменение непредсказуемо распространяется по ячейке. Поскольку тег участвует в обработке следующей ячейки, последующие ячейки этой цепи также становятся нерасшифровываемыми, и цепь приходится отбросить. Поэтому атакующий может лишь уничтожить цепь, что он и без того способен сделать.
\end{soluzione}

\begin{soluzione}{es:10-padding}
Наблюдаемые интервалы: 4,1; 3,0; 5,7; 8,7 и 6,8 секунды, среднее 5,7. Для максимума двух независимых равномерных величин на $[a, b]$ математическое ожидание равно $a + \tfrac{2}{3}(b - a) = 1.5 + \tfrac{2}{3} \cdot 8 \approx 6.8$ секунды. Пять наблюдений — слишком малая выборка, чтобы разница была удивительной. Кроме того, таймер перезапускается при прохождении через канал \emph{любой} ячейки, а не только ячеек заполнения.
\end{soluzione}

\begin{soluzione}{es:10-drop}
В лабораторной сети промежуточный ретранслятор не отвечает на \texttt{DROP}; сообщение распознаётся и молча отбрасывается. Косвенное доказательство — цепь остаётся пригодной: после двух сообщений \texttt{DROP} клиент всё ещё может расширить её до выходного ретранслятора. Если бы промежуточный ретранслятор не распознал ячейки, будучи в тот момент последним переходом цепи, он уничтожил бы цепь. Если бы клиент и промежуточный ретранслятор по-разному учитывали ячейки, их накопительные дайджесты рассинхронизировались бы и последующее расширение завершилось бы неудачей.
\end{soluzione}

\begin{soluzione}{es:10-alterazione}
В тесте приватной сети изменение одного бита в ячейке \texttt{EXTEND2}, предназначенной входному узлу, который в тот момент является единственным переходом цепи, приводит к неудаче проверки дайджеста. Как последний узел этой цепи входной ретранслятор отвечает ячейкой \cella{DESTROY} с причиной 1 (\texttt{PROTOCOL}). В более длинной цепи изменение, внесённое клиентом в ячейку для третьего перехода, было бы переслано через первые два перехода и обнаружено только третьим.
\end{soluzione}
